%%%%%%%%%%%%%%%%%% 选择公理、良序定理与佐恩引理 %%%%%%%%%%%%%%%%%%%

\chapter{选择公理、良序定理与佐恩引理}

选择公理对于大部分数学家来说在直觉上都是正确的，不过却可以用它推导出一些令人觉得深奥、较不合乎直觉的定理或性质，例如良序定理和佐恩引理（事实上这三个性质是等价的）。对于选择公理的使用，有些数学家并不一定那么释怀，不过由于过去一些数学证明常不自知地用到选择公理，而且这个公理的使用并没有和数学理论造成任何矛盾，所以目前绝大部分的数学家是接受选择公理以及与其等价的良序定理和佐恩引理的。也因此将来大家在学习进阶的一些数学课程时会遇到需要用这三个性质其中之一所得到的定理，所以我们特别介绍这三个性质。要注意，我们并不是要证明这三个性质是等价的（有兴趣的同学可参考集合论相关的书籍），而是着重于了解这三个性质以及如何运用。

\section{选择公理}

选择公理指的是对于任意的非空集合 $S$，我们都可以定义一个 “方法” 在 $S$ 的任意非空子集中挑出一个代表元素。这个性质在 $S$ 是有限集时不会有问题，因为此时 $S$ 的非空子集只有有限多个，我们可以用列举的方法告诉大家如何挑选每个非空子集的代表元素。另外对于一些特别的无限集，也可能没问题。例如当 $S$ 为 $\mathbb{N}$ 时，我们可以对于 $\mathbb{N}$ 的每个非空子集利用良序原理，使用选最小元的方法挑出代表元素。不过对于更一般的无限集，就可能会有问题了，因为我们可能无法明确地造出 “挑选” 的方法。不过换一个角度来说，既然是非空子集，我们可以任意挑选一个元素当代表元素啊！所以虽然无法像有限集的情况用列举法，不过至少应该还是一种 “挑选” 的方法。因此一般的数学家认为这个性质只是对于有限集的情形延伸而来，颇符合直觉，也因此公认它是对的。

要注意，前述在非空子集中任意挑选一个元素当代表元素的说法其实是有限制的。既然要挑代表元素，便要有一致性。也就是说对于同一个非空子集，不能一下子挑一个元素，一下子又挑另一个元素。这种一致性利用函数来表达便最合适了。回顾一下集合 $S$ 的子集所成的集合，即所谓 $S$ 的幂集，我们用 $\mathcal{P}(S)$ 来表示。我们想将 $\mathcal{P}(S)$ 中任一个非空子集 $A$ 对应到 $A$ 的某一个元素，所以用函数的观点来看就是找到一个函数 $f$ 使得对任意 $S$ 的非空子集 $A$ 皆有 $f(A) \in A$。这一个函数 $f$，由于是挑选 $S$ 中每一个非空子集的代表元素，所以一般我们也称之为选择函数（choice function）。接下来，我们写下选择公理的正确叙述。

\textbf{选择公理（Axiom of Choice）:} {\kaishu 对于任何非空集合 $S$，存在一个选择函数 $f : \mathcal{P}(S) \setminus \{\emptyset\} \to S$，使得对于 $S$ 的每一个非空子集 $A$（即 $A \in \mathcal{P}(S) \setminus \{\emptyset\}$），我们有 $f(A) \in A$。}\index{选择公理（Axiom of Choice）}

要注意，这里 $A$ 虽然是一个集合，但选择函数对其作用后所得的是 $A$ 中的一个元素，而不是集合。也就是说选择函数 $f$ 的定义域是 $S$ 的非空子集合，而不是 $S$，千万不要误以为 $f(A)$ 是前一章 5.2 节所提的 “$A$ 在 $f$ 下的像”。另外要说明的是选择公理中仅谈到选择函数的存在性，并未论及如何找到此选择函数。所以一般利用选择公理所得的结果，它的存在性都不会是构造性的。

前面提及，我们常不经意地用到选择公理。例如在命题 5.5.9 的证明中，我们事实上用到了选择公理。也就是在 $S_1, \dots, S_n, \dots$ 皆为可数集的假设上，对于每一个 $i \in \mathbb{N}$，由于 $|S_i| \leq |\mathbb{N}|$，我们就在可能很多 $S_i \to \mathbb{N}$ 的单射函数中挑选了一个 $f_i : S_i \to \mathbb{N}$ 作为代表。所以严格来说命题 5.5.9 是要用到选择公理才能成立的。最后，我们再看一个利用选择公理所得的结果。

\begin{proposition}
假设 $S$ 为无限集，则 $S$ 中存在子集是可数无限的。
\end{proposition}

\begin{proof}
假设 $f : \mathcal{P}(S) \setminus \{\emptyset\} \to S$ 为选择函数。考虑 $g : \mathbb{N} \to S$，定义为 $g(1) = f(S)$。令 $S_2 = S \setminus \{f(S)\}$，由于 $S$ 为无限集，我们有 $S_2 \neq \emptyset$，也就是说 $S_2 \in \mathcal{P}(S) \setminus \{\emptyset\}$，所以 $f(S_2)$ 是有定义的，现令 $g(2) = f(S_2)$。现利用数学归纳法对于 $k \geq 2$ 令 $S_{k+1} = S \setminus \{f(S), f(S_2), \dots, f(S_k)\}$。由 $S$ 为无限集，我们知 $S_{k+1} \in \mathcal{P}(S) \setminus \{\emptyset\}$，故定义 $g(k+1) = f(S_{k+1})$。依此，我们定义了一个 $g : \mathbb{N} \to S$ 这样一个单射函数，也因此 $g$ 的像，即 $g(\mathbb{N})$ 就是 $S$ 的一个可数无限的子集。
\end{proof}

一般来说，在处理有无穷多个集合的问题时，就有可能用到选择公理，但也未必一定要用到。英国哲学家罗素（Russell，也是数学家和逻辑学家）提出了一个简单区分是否要用到选择公理的例子。当在无穷多双的袜子里（一般来说每双袜子都是同色且不分左右脚），若要从每双袜子中都挑出一只袜子，便需用到选择公理；但在无穷多双的鞋子里（一般来说每双鞋子都有分左右脚），若要从每双鞋子里都挑出一只鞋子，便不需用到选择公理，因为我们可以简单地都挑左脚鞋。简单来说，是否用到选择公理取决于函数是否为构造性的。

\section{良序定理}

回顾一下，一个偏序集 $(X, \preceq)$ 为全序的表示 $X$ 中任意两个元素都可以比较大小，亦即若 $a, b \in X$，则 $a \preceq b$ 或 $b \preceq a$ 其中之一必成立。而一个偏序集 $(X, \preceq)$ 称为良序的，表示每一个 $X$ 的非空子集 $T$ 都会有最小元，也就是说存在 $t_0 \in T$ 满足 $t_0 \preceq t, \forall t \in T$。为了方便起见我们用 $\min(T)$ 表示 $T$ 的最小元。所谓良序定理（well-ordering theorem），就是说对于一个非空的集合 $X$，我们都可以找到一个序 $\preceq$，使得 $(X, \preceq)$ 是一个良序集。

不要将良序定理和良序原理混淆。良序原理简单来说指的是用我们一般的序 $\leq$ 会使得 $(\mathbb{N}, \leq)$ 是良序集。而良序定理指的是任意的非空集合 $X$，另外它并未指出哪一种序 $\preceq$ 会使得 $(X, \preceq)$ 是一个良序集。例如有理数 $\mathbb{Q}$ 利用一般的 $\leq$，并不会是良序的（我们无法在 $\{r \in \mathbb{Q} \mid 0 < r < 1\}$ 这个集合中找到最小的有理数）。不过利用 $\mathbb{Q}$ 是可数的（推论 5.5.8），我们可以利用 $\mathbb{N}$ 和 $\mathbb{Q}$ 的一一对应，将 $\mathbb{Q}$ 重新排序，而利用此新的排序得到 $\mathbb{Q}$ 为良序的。所以说良序定理对于可数集，我们知道可由 $\mathbb{N}$ 的良序原理推得。不过对于不可数集，良序定理就较难令人理解。事实上到目前为止，我们仍无法具体举出实数 $\mathbb{R}$ 上的序 $\preceq$ 使得 $(\mathbb{R}, \preceq)$ 是良序的。

良序定理又称为策梅洛定理（Zermelo's theorem），\index{策梅洛定理（Zermelo's theorem）|see{良序定理（well-ordering theorem）}}它是策梅洛利用选择公理证出的。我们可以很快地利用良序定理证出选择公理。所以逻辑上，它们是等价的。因此我们也可以将良序定理视为是一种公理。不过如前所述，一般人直觉上较能觉得选择公理是对的，而直觉上较无法理解良序定理，所以一般我们不称它为公理，觉得它是由选择公理所推出的定理。我们正式写下其定理形式。

\begin{theorem}[良序定理]\index{良序定理（well-ordering theorem）}
令 $X$ 为一个非空集合。则在 $X$ 上存在一个序关系 $\preceq$ 使得 $(X, \preceq)$ 是良序的。
\end{theorem}

良序定理强调的是可以找到序 $\preceq$，使得 $(X, \preceq)$ 为良序的。不过由于这个存在性是利用选择公理所得到，前面强调过用选择公理所得的存在性不是构造性的，所以良序定理并无法提出使得 $(X, \preceq)$ 为良序集的序 $\preceq$ 为何。

在这里我们不去探讨如何由选择公理得到良序定理。不过反向是容易的，即利用良序定理推得选择公理。对于任意的非空集合 $S$，利用良序定理，考虑序 $\preceq$ 使得 $(S, \preceq)$ 为良序的。此时对任意 $A \in \mathcal{P}(S) \setminus \{\emptyset\}$，令 $f(A) = \min(A)$。则 $f : \mathcal{P}(S) \setminus \{\emptyset\} \to S$ 满足选择函数的要求，即 $f(A) \in A$。故推得选择公理。

有了良序定理这个强大的工具，就如同 $\mathbb{N}$ 有良序原理一样，我们可以有类似数学归纳法的方法，称之为超限归纳法（transfinite induction）。回顾一下推论 2.3.6，要使用数学归纳法证明 $P(n)$ 对所有 $n \in \mathbb{N}$ 皆成立，我们先证明
\begin{enumerate}
    \item $P(1)$ 是对的；
    \item 若对任意 $i < k$, $P(i)$ 皆成立，则 $P(k)$ 亦成立。
\end{enumerate}
证明了 (1)、(2) 便证得 $P(n)$ 对所有 $n \in \mathbb{N}$ 皆成立。而超限归纳法是利用 $(X, \preceq)$ 为良序的，先证明
\begin{enumerate}
    \item $P(\min(X))$ 成立；
    \item 假设 $\beta \in X$。若对任意 $\alpha \prec \beta$, $P(\alpha)$ 皆成立，则 $P(\beta)$ 亦成立。
\end{enumerate}
如此便证得了 $P(x)$ 对所有 $x \in X$ 皆成立。我们将此定理叙述如下：

\begin{theorem}[超限归纳法]\index{超限归纳法（transfinite induction）}
假设 $(X, \preceq)$ 为良序集且令 $x_1 = \min(X)$。假设以下两个命题是对的，那么对任意 $x \in X$, $P(x)$ 皆会成立。
\begin{enumerate}
    \item $P(x_1)$ 成立。
    \item 假设 $\beta \in X$。若对任意 $\alpha \in X$ 满足 $\alpha \prec \beta$, $P(\alpha)$ 皆成立，则 $P(\beta)$ 成立。
\end{enumerate}
\end{theorem}

\begin{proof}
利用反证法，假设存在 $x \in X$ 使得 $P(x)$ 不成立。考虑 $S = \{x \in X \mid P(x) \text{ 不成立}\}$。依假设 $S \neq \emptyset$。故由 $(X, \preceq)$ 为良序的，知存在 $\beta \in S$ 使得 $\beta = \min(S)$。由 (1) 知 $\beta \neq x_1$。又因 $\beta = \min(S)$，我们知对任意 $\alpha \in X$ 满足 $\alpha \prec \beta$，皆有 $\alpha \notin S$，亦即 $P(\alpha)$ 皆成立。故由 (2) 知 $P(\beta)$ 成立。此与 $\beta \in S$ 之假设相矛盾，故知 $S = \emptyset$，亦即对任意 $x \in X, P(x)$ 皆成立。
\end{proof}

良序定理的好处就是可以让我们如同处理 $\mathbb{N}$ 一样地处理一般的集合。不过这反而会造成同学的误解。许多同学会误以为一个集合是良序的，表示可以由最小的元素开始排序。若将最小的元素对应到 1，第二小的元素对应到 2，这样一直下去不就表示所有无限集都可以和 $\mathbb{N}$ 形成一一对应，而造成所有的集合都是可数的这样的奇怪现象？这样的看法其实是错的，主要的原因是我们确实可以利用良序的性质将集合中的元素从小排到大，但这并不代表可以将每一个元素都排到。例如对于 $\mathbb{N}$ 我们可以有以下的序 $\preceq$：若 $a, b \in \mathbb{N}$ 且同奇同偶，则定 $a \preceq b$ 当且仅当 $a \leq b$；而若 $a$ 为奇数 $b$ 为偶数，则定 $a \preceq b$。在此定义之下我们仍可得 $(\mathbb{N}, \preceq)$ 为良序集，但是若从小到大排序，则 2（或任何的偶数）永远都无法被排到（因为必须将奇数先排完才能排偶数）。

\section{佐恩引理}

佐恩引理（Zorn's lemma）和选择公理也是等价的。由于它的前提对一般的偏序集都可使用，所以较易拿来应用。也因此有许多数学上的定理是用它推得的，故一般称之为引理。事实上，以后大家在遇到符合佐恩引理的前提时，可以如使用公理一样直接套用。所以我们不去论证佐恩引理和选择公理（或良序定理）之间的等价关系，而专注于了解这个引理。

首先我们再回顾一些有关于序的定义。在一个偏序集 $(S, \preceq)$ 中，假设 $T$ 是 $S$ 的子集且 $T$ 在 $\preceq$ 之下 $(T, \preceq)$ 是一个全序集（意即对任意 $t, t' \in T$ 皆有 $t \preceq t'$ 或 $t' \preceq t$），我们称 $T$ 为 $(S, \preceq)$ 的一个链（chain）。对于 $S$ 的一个非空子集 $S_0$，我们说 $u \in S$ 是 $S_0$ 的一个上界，表示对任意 $s_0 \in S_0$ 皆有 $s_0 \preceq u$。而我们说 $\mu \in S$ 是 $S$ 的极大元（maximal element），表示 $\mu \in S$ 且 $S$ 中没有任何元素 $s$ 会满足 $\mu \prec s$（也就是说 $S$ 中的元素 $s$，要不是满足 $s \preceq \mu$ 就是和 $\mu$ 不能比较大小）。现在我们可以叙述佐恩引理。

\begin{lemma}[佐恩引理]
假设 $(S, \preceq)$ 是一个偏序集。若 $S$ 中每一个链皆在 $S$ 中有上界，则 $S$ 有极大元。
\end{lemma}

从佐恩引理的叙述可知，当我们要证明一个特定的偏序集有极大元，可以考虑使用佐恩引理。也就是说，若要证明一个偏序集 $(S, \preceq)$ 有极大元，我们可以考虑 $S$ 中任意的一组链，然后试着在 $S$ 中找到此组链的上界。如果能证明所有的链皆在 $S$ 中可找到上界，那么偏序集 $(S, \preceq)$ 会有极大元。要注意，这里我们特别强调，每一个链的上界必须在 $S$ 中找到，否则不能确保 $S$ 有极大元。

可以看出佐恩引理仍然是谈论存在性的问题。由于它和选择公理以及良序定理是等价的，所以它所得的存在性也不是构造性的。也就是说，我们可证得极大元的存在性，但无从得知这些极大元有哪些。佐恩引理会在将来许多数学课程中用到。例如代数中，要证明每个环（ring）皆存在着极大理想（maximal ideal）；线性代数中，要证明所有的（无限维）向量空间（vector space）皆存在一组基（basis），都要用到佐恩引理。这里我们特别举出一个常常使用佐恩引理的情况，事实上前面所举的这两类例子就是在这种情况之下得证的。

\begin{proposition}
假设 $\mathcal{S}$ 是由一些具有某特定性质的集合（sets）所成的集合。考虑一般集合的包含关系 $\subseteq$ 所形成的偏序集 $(\mathcal{S}, \subseteq)$。若 $\mathcal{S}$ 中任意的一组链的并集仍在 $\mathcal{S}$ 中，则在 $\mathcal{S}$ 中必存在一集合 $M$ 使得 $\mathcal{S}$ 中任一集合 $S$ 皆不会满足 $M \subsetneq S$。
\end{proposition}

\begin{proof}
注意 $\mathcal{S}$ 这是一个集合里的元素仍为集合。很明显地 $\mathcal{S}$ 中任一组链的并集必包含此链中任一集合，故由 $\mathcal{S}$ 中任意的一组链的并集仍在 $\mathcal{S}$ 中之假设知此并集便是此链的一个上界。因此由佐恩引理知 $(\mathcal{S}, \subseteq)$ 有极大元 $M$，亦即 $M$ 是 $\mathcal{S}$ 中的一个集合，而且 $\mathcal{S}$ 中没有其他的集合会比 $M$ “大”（也就是说没有其他的集合会包含 $M$），故得证本定理。
\end{proof}

这里我们要特别说明的是，任意的一组集合的并集当然包含其中任一集合，所以或许同学会疑问为何命题 6.3.2 中要假设任意的一组链的并集仍在 $\mathcal{S}$ 中？这就是我们前面强调的要利用佐恩引理的重点在于每一个链的上界要在 $\mathcal{S}$ 中。所以光取并集虽可符合上界的要求，但可能不符合在 $\mathcal{S}$ 中的要求。因此一般使用命题 6.3.2 的重点在于如何利用 “链” 的特性，证明将它们取并集后仍会在 $\mathcal{S}$ 中。这一点将来遇到用佐恩引理证明问题时，务必留意。